
\documentclass{article}
\usepackage{colortbl}
\usepackage{makecell}
\usepackage{multirow}
\usepackage{supertabular}

\begin{document}

\newcounter{utterance}

\centering \large Interaction Transcript for game `cladder', experiment `full\_v1.5\_default', episode 5310 with refineonly{-}Qwen3.5{-}27B.
\vspace{24pt}

{ \footnotesize  \setcounter{utterance}{1}
\setlength{\tabcolsep}{0pt}
\begin{supertabular}{c@{$\;$}|p{.15\linewidth}@{}p{.15\linewidth}p{.15\linewidth}p{.15\linewidth}p{.15\linewidth}p{.15\linewidth}}
    \# & \multicolumn{2}{c}{Player} && \multicolumn{2}{c}{Game Master} \\
    \hline

    \theutterance \stepcounter{utterance}  
    & & & \multicolumn{4}{p{0.6\linewidth}}{
        \cellcolor[rgb]{0.9,0.9,0.9}{
            \makecell[{{p{\linewidth}}}]{
                \texttt{\tiny{[P1$\langle$GM]}}
                \texttt{You are an expert in causal inference. The following question is not a typical commonsense query, but rather a meticulously designed question created by a professor specializing in causal inference, intended to assess the students' mastery of the course content.} \\
\texttt{Imagine a self{-}contained, hypothetical world with only the following conditions, and without any unmentioned factors or causal relationships: Liking spicy food has a direct effect on occupation and salary. Skill has a direct effect on occupation and salary. Occupation has a direct effect on salary. Skill is unobserved. Method 1: We look at how liking spicy food correlates with salary case by case according to occupation. Method 2: We look directly at how liking spicy food correlates with salary in general. To understand how liking spicy food affects salary, is it more correct to use the Method 1 than Method 2?} \\
\texttt{Start your answer with 'Yes' or 'No', followed by additional reasoning or evidence to support your explanation.} \\
            }
        }
    }
     \\ \\

    \theutterance \stepcounter{utterance}  
    & \multicolumn{4}{p{0.6\linewidth}}{
        \cellcolor[rgb]{0.561,0.737,0.561}{
            \makecell[{{p{\linewidth}}}]{
                \texttt{\tiny{[P1$\rangle$GM]}}
                \texttt{No} \\
\\ 
\texttt{In this hypothetical world, **Skill** acts as an **unobserved confounder** because it causally influences both the treatment (Liking spicy food) and the outcome (Salary). This creates a backdoor path: Liking spicy food \$\textbackslash{}leftarrow\$ Skill \$\textbackslash{}rightarrow\$ Salary.} \\
\\ 
\texttt{1.  **Method 2 (General Correlation)**: This estimates the total association between Liking spicy food and Salary. Because Skill is unobserved, this estimate is biased by the confounding effect of Skill. It captures the true causal effect plus the spurious correlation induced by Skill.} \\
\texttt{2.  **Method 1 (Stratification by Occupation)**: This method conditions on **Occupation**. In the specified causal graph, Occupation is a **mediator** (Liking spicy food \$\textbackslash{}rightarrow\$ Occupation \$\textbackslash{}rightarrow\$ Salary) and also a descendant of the unobserved confounder (Skill \$\textbackslash{}rightarrow\$ Occupation).} \\
\texttt{    *   **Blocking the Causal Path**: By conditioning on a mediator, Method 1 blocks the indirect causal pathway through which Liking spicy food affects Salary. Consequently, it fails to estimate the *total* effect, capturing only the direct effect (if identifiable).} \\
\texttt{    *   **Failure to Remove Confounding**: Crucially, conditioning on Occupation does not block the backdoor path through Skill. Since Skill affects both Liking spicy food and Occupation, and we are conditioning on Occupation, the association between Liking spicy food and Salary within each stratum remains confounded by Skill. In some scenarios, conditioning on a descendant of a confounder can even induce or exacerbate bias (collider bias), though the primary issue here is the persistence of the unblocked backdoor path.} \\
\\ 
\texttt{Therefore, Method 1 is **not more correct** than Method 2. Method 1 introduces bias by blocking the causal mechanism (mediation) and fails to eliminate the confounding bias from the unobserved Skill. Since Skill is unobserved, the total causal effect is unidentifiable using either method alone, but Method 1 is structurally flawed for estimating the total effect because it alters the causal quantity of interest and does not resolve the confounding.} \\
            }
        }
    }
    & & \\ \\

    \theutterance \stepcounter{utterance}  
    & & & \multicolumn{2}{p{0.3\linewidth}}{
        \cellcolor[rgb]{0.9,0.9,0.9}{
            \makecell[{{p{\linewidth}}}]{
                \texttt{\tiny{[GM$|$GM]}}
                \texttt{No} \\
            }
        }
    }
    & & \\ \\

    \theutterance \stepcounter{utterance}  
    & & & \multicolumn{2}{p{0.3\linewidth}}{
        \cellcolor[rgb]{0.9,0.9,0.9}{
            \makecell[{{p{\linewidth}}}]{
                \texttt{\tiny{[GM$|$GM]}}
                \texttt{no} \\
            }
        }
    }
    & & \\ \\

    \theutterance \stepcounter{utterance}  
    & & & \multicolumn{2}{p{0.3\linewidth}}{
        \cellcolor[rgb]{0.9,0.9,0.9}{
            \makecell[{{p{\linewidth}}}]{
                \texttt{\tiny{[GM$|$GM]}}
                \texttt{game\_result = WIN} \\
            }
        }
    }
    & & \\ \\

\end{supertabular}
}

\end{document}
